% !TeX root = ../../Book1.tex
% 纲要标题：### C.1 自由积与带合并自由积
% 纲要位置：计划纲要.md，第 889 行。
% BEGIN APPENDIX OUTLINE SYNC
% #### C.1.1 自由积的约化字
% #### C.1.2 识别公共子群的构造
% #### C.1.3 标准形与因子的嵌入
% END APPENDIX OUTLINE SYNC

\section{自由积与带合并自由积}
\label{b1:sec:C01}

把两个群放进同一个群中，可以要求它们的元素逐一交换，也可以只保留各自已有的关系。直积与自由积分别对应这两种要求。若两个群还包含一个需要认作相同的子群，就应在自由积中加入识别关系；关键是确认这些关系不会把同一因子中原本不同的元素意外识别。

\subsection{自由积的约化字}
\label{b1:subsec:C01:01}

\cref{b1:lem:free-product-normal-form}已经证明：自由积 \(A*B\) 的每个元素有唯一交替字 \(x_1\cdots x_n\)，各字母来自 \(A\setminus\{1\}\) 或 \(B\setminus\{1\}\)，相邻字母来自不同因子；空字表示单位。这里一个字母可以是因子中的任意非单位元素，不限于某组生成元。化简时把相邻的同因子字母相乘，并删除所得的单位元，直到非单位字母交替为止。两因子的同态可以由\cref{b1:thm:free-product-universal}唯一拼成自由积上的同态，不要求它们的像交换。

\ProseParagraphBreak
例如 \(C_2*C_3=\langle a,b\mid a^2=b^3=1\rangle\)。在字 \((ab)^m\) 中，\(a,b\) 均非单位元，连续字母始终来自不同因子，没有可以相乘化简的相邻字母。因此对每个 \(m>0\)，这个字都约化且非空，由标准形知 \((ab)^m\ne1\)，所以 \(ab\) 有无限阶。两个有限群的自由积可以无限；这与直积 \(C_2\times C_3\cong C_6\) 的情形不同。这个判断来自标准形，不能只从展示中“看不见 \((ab)^m=1\)”便推断其不成立。

\subsection{识别公共子群的构造}
\label{b1:subsec:C01:02}

所谓“公共子群”，可以是同一个抽象群在两个因子中的不同实现，无须原先就是两因子的同一个集合子群。例如取两个无限循环群 \(\langle a\rangle\)、\(\langle b\rangle\)，希望将第一个群中的 \(a^2\) 与第二个群中的 \(b^3\) 认作同一元素，同时保留两个因子中的其他运算。这时取抽象无限循环群 \(C=\langle c\rangle\)，指定嵌入 \(i(c)=a^2\)、\(j(c)=b^3\)，再加入关系 \(a^2=b^3\)。一般构造用两条指定的单同态记录这种要求，识别的是同一个 \(c\in C\) 在两因子中的像 \(i(c),j(c)\)。

\begin{definition}\label{b1:def:amalgamated-free-product}
设 \(C,A,B\) 为群，\(i:C\hookrightarrow A\)、\(j:C\hookrightarrow B\) 为单同态。把 \(A,B\) 视为自由积中的两个因子，并令
\[
N=\bigl\langle\!\bigl\langle i(c)j(c)^{-1}\mid c\in C\bigr\rangle\!\bigr\rangle
\trianglelefteq A*B.
\]
商群 \((A*B)/N\) 称为沿 \(C\) 的\term[dai-hebing-ziyouji]{带合并自由积}[free product with amalgamation]，记为 \(A*_C B\)。记从两因子到该商群的自然同态为 \(u:A\rightarrow A*_C B\)、\(v:B\rightarrow A*_C B\)。
\end{definition}

\symidx{\(A*_C B\)：沿给定嵌入合并 \(C\) 的自由积}

这里的双尖括号表示正规闭包，即包含所列元素的最小正规子群。若商映射记为 \(\pi\)，则 \(u(a)=\pi(a)\)、\(v(b)=\pi(b)\)；关系 \(i(c)j(c)^{-1}\in N\) 恰使 \(u\bigl(i(c)\bigr)=v\bigl(j(c)\bigr)\)。一旦令这个关系元素在商群中等于 \(1\)，它的每个共轭 \(w\,i(c)j(c)^{-1}w^{-1}\)（\(w\in A*B\)）以及这些共轭及其逆的有限乘积也必须等于 \(1\)。这正是\cref{b1:def:normal-closure}及其后的有限乘积描述中，关系通过共轭与乘积传播的规则。取正规闭包既保证商群所用的核正规，也只加入这些识别所强制的关系。

\ProseParagraphBreak
下标 \(C\) 还隐含了两条嵌入 \(i,j\)；只知道三个抽象群，并不足以指定这个构造。定义立即给出 \(ui=vj\)，但加入关系的商构造本身并不排除因子中的非单位元素被送到单位元，所以 \(u,v\) 的单射性还需要下一小节的证明。对于任一群 \(H\) 及同态 \(f:A\to H\)、\(g:B\to H\)，条件 \(fi=gj\) 表示：对每个 \(c\in C\)，两条路径给出同一个像 \(f(i(c))=g(j(c))\)。泛性质断言，这样的两条同态能由商群上的唯一同态同时得到。

\begin{proposition}\label{b1:prop:amalgam-universal}
设 \(C,A,B,H\) 为群，\(i:C\hookrightarrow A\)、\(j:C\hookrightarrow B\) 为单同态，\(A*_C B\) 是沿 \(i,j\) 构造的带合并自由积，\(u:A\rightarrow A*_C B\)、\(v:B\rightarrow A*_C B\) 为自然同态。若同态 \(f:A\rightarrow H\)、\(g:B\rightarrow H\) 满足 \(fi=gj\)，则存在唯一同态 \(h:A*_C B\rightarrow H\) 满足 \(hu=f\)、\(hv=g\)。
\end{proposition}

\begin{proof}
自由积的泛性质给出唯一同态 \(\widetilde h:A*B\rightarrow H\)，在两因子上的限制分别是 \(f,g\)。对每个 \(c\in C\)，有
\[
\widetilde h\bigl(i(c)j(c)^{-1}\bigr)
=f\bigl(i(c)\bigr)g\bigl(j(c)\bigr)^{-1}=1.
\]
核是正规子群，故包含定义中的 \(N\)。由\cref{b1:thm:quotient-universal}，\(\widetilde h\) 唯一地通过商群下降为 \(h\)。两因子的像生成商群，因此任何满足所述限制的同态都只能是这个 \(h\)。
\end{proof}

当 \(C=1\) 时，没有额外识别，得到普通自由积。若 \(C=A=B\) 且两嵌入均为恒等，两个因子的对应元素全被识别，泛性质给出 \(A*_A A\cong A\)。

\subsection{标准形与因子的嵌入}
\label{b1:subsec:C01:03}

利用 \(i,j\) 可把 \(C\) 分别视为 \(A,B\) 的子群。为每个右陪集 \(Ca\)、\(Cb\) 选一个代表，得到代表集 \(T_A,T_B\)，均以 \(1\) 代表 \(C\) 本身。若 \(t\) 代表 \(Ca\)，则 \(Ca=Ct\)，从而 \(a=ct\)；使用这种右陪集约定，\(C\) 系数位于代表的左边。因此每个因子元素都有唯一分解 \(ct\)，其中 \(c\in C\)，\(t\) 属于该因子的代表集。这里对非空陪集族应用\cref{b1:def:A-choice}中的选择公理，再把单位陪集的代表改取为 \(1\)。这保证代表集存在，却没有给出从输入元素求出 \(c,t\) 的算法；有效计算还需另有可执行的陪集分解方法。

\ProseParagraphBreak
这一步的唯一性来自两个不同层次。给定 \(a\in A\)，陪集 \(Ca\) 先确定唯一代表 \(t\in T_A\)，再由 \(a=ct\) 确定 \(c=at^{-1}\in C\)。不同因子的代表仍保留各自的标记；它们只是各因子元素的一种写法。我们要进一步证明，在合并所得的群中，也可以用这些代表写出唯一的交替表达式。由于 \(C\) 未必在因子的中心，不能简单地把一个乘积中的所有 \(C\) 元素交换到最左边；下述证明会在每次左乘时重新作陪集分解。

\begin{theorem}[带合并自由积的标准形]\label{b1:thm:amalgam-normal-form}
设 \(C,A,B\) 为群，\(i:C\hookrightarrow A\)、\(j:C\hookrightarrow B\) 为单同态，并通过 \(i,j\) 把 \(C\) 分别视为 \(A,B\) 的子群。分别选取 \(C\) 在 \(A,B\) 中的右陪集代表集 \(T_A,T_B\)，并以 \(1\) 代表单位陪集。则每个 \(A*_C B\) 中的元素唯一表示为
\[
c t_1\cdots t_n,
\qquad c\in C,\quad t_k\in(T_A\setminus\{1\})\sqcup(T_B\setminus\{1\}),
\]
其中各 \(t_k\) 交替来自两代表集，允许 \(n=0\)。特别地，自然映射 \(A\rightarrow A*_C B\)、\(B\rightarrow A*_C B\) 均单射，其像的交恰为 \(C\) 的像。
\end{theorem}

\begin{proof}
只证明每个字都能化成这种形式，还不足以排除两个不同的标准串在商群中相等。为证明唯一性，我们构造一个群作用，使每个候选标准串都把同一个基点送到该串所标记的位置。不同串给出不同位置，而同一群元素在基点上的作用只能有一个值，因而可以用这个作用区分标准串。

令 \(\Omega\) 为所有满足定理条件的有限有序串 \((c;t_1,\ldots,t_n)\) 组成的集合。此时相等只表示长度相同、各位置的符号相同，尚未使用商群中的相等关系。没有代表项且 \(c=1\) 的串记为 \(\omega_0=(1;)\)。对 \(a\in A\)，定义映射 \(L_a:\Omega\rightarrow\Omega\)。若首个代表 \(t_1\) 来自 \(T_A\)，则 \(a,c,t_1\) 全在 \(A\) 中，左乘 \(a\) 后可以直接合并这个前缀，代表 \(t_1\) 也可能被消去。在 \(A\) 中把 \(act_1\) 唯一分解为 \(c't'\)，并令
\[
L_a(c;t_1,t_2,\ldots,t_n)=(c';t',t_2,\ldots,t_n).
\]
若原串没有代表项，或 \(t_1\) 来自 \(T_B\)，则先把 \(ac\) 分解为 \(c't'\)，再把新得到的 \(A\) 代表插在原代表串之前，令
\[
L_a(c;t_1,\ldots,t_n)=(c';t',t_1,\ldots,t_n).
\]
这两个公式中，若 \(t'=1\)，均删去该代表项。所得串仍交替：第一种情形中余下的首项 \(t_2\) 若存在，便属于另一因子；第二种情形中原来的 \(t_1\) 若存在，也属于另一因子。对 \(b\in B\) 交换两因子的角色，定义 \(L_b\)。

对同一因子中的 \(a,a'\)，有 \(L_aL_{a'}=L_{aa'}\)，其中左边先作用 \(L_{a'}\)。这个等式取决于因子中 \(ct\)-分解的唯一性。若原串以该因子的代表开头，设第一次分解为 \(a'ct_1=c't'\)。当 \(t'\ne1\) 时，第二次分解 \(ac't'=aa'ct_1\)；当 \(t'=1\) 时，原首代表已消去，第二次分解 \(ac'=aa'ct_1\)，余下的代表串从另一因子开始。两种情形的前缀都是同一个元素 \(aa'ct_1\)，所以与一次作用 \(L_{aa'}\) 得到的分解相同。若原串不以该因子开头，只需将这里的 \(ct_1\) 换成 \(c\)，两次操作最终分解的就是 \(aa'c\)。又 \(L_1\) 是恒等映射，故 \(L_{a^{-1}}\) 是 \(L_a\) 的逆，得到同态 \(A\rightarrow S_\Omega\)，其中 \(S_\Omega\) 是集合 \(\Omega\) 上的全体自双射所成的群；\(B\) 亦然。

对于 \(c_0\in C\)，若首代表来自正在使用的因子，则 \(c_0ct_1\) 已有分解 \((c_0c)t_1\)；若首代表来自另一因子或不存在，则 \(c_0c\in C\) 的代表为 \(1\)。因此两作用都只把首项 \(c\) 改成 \(c_0c\)，在公共子群上相容。以 \(H=S_\Omega\) 应用\cref{b1:prop:amalgam-universal}，遂得到 \(A*_C B\) 在 \(\Omega\) 上的作用。

任意因子字从右向左依次作用于 \(\omega_0\)，每一步的前缀分解都在商群中保持元素不变，所以输出串表示原元素。这证明存在性。另一方面，标准串对应的群元素 \(ct_1\cdots t_n\) 作用于 \(\omega_0\) 时，先加入 \(t_n\)，再加入 \(t_{n-1}\)，依次进行；因代表交替，过程中没有同因子合并。最后左乘 \(c\)，得到
\[
(ct_1\cdots t_n)\cdot\omega_0=(c;t_1,\ldots,t_n).
\]
如果两个标准串表示同一群元素，它们在已经构造好的群作用中必给出同一个置换，在 \(\omega_0\) 上的取值也相同，故串完全相同。这证明唯一性。

因子中的元素 \(ct\) 只有在 \(c=t=1\) 时才等于单位，故两自然映射单射。若一个 \(A\) 中的元素等于一个 \(B\) 中的元素，标准串唯一性迫使两者的非单位代表都不出现，因而均属于 \(C\)。
\end{proof}

例如取 \(A=\langle a\rangle\cong\mathbb Z\)、\(B=\langle b\rangle\cong\mathbb Z\)，把 \(C\cong\mathbb Z\) 的生成元分别送到 \(a^2,b^3\)。则
\[
A*_C B\cong\langle a,b\mid a^2=b^3\rangle.
\]
可取 \(T_A=\{1,a\}\)、\(T_B=\{1,b,b^2\}\)。记公共元素 \(z=a^2=b^3\)，标准形是 \(z^m t_1\cdots t_n\)。这既说明两个无限循环因子没有被压缩，也给出检验元素是否相等的方法；其中 \(z\) 与 \(a,b\) 都交换，故属于整个群的中心。

\ProseParagraphBreak
例如左乘 \(a\) 于标准形 \(z^mab^2\)，首两个 \(a\) 合并为 \(z\)，结果为 \(z^{m+1}b^2\)；左乘 \(b\) 于 \(z^ma\)，则得到 \(z^mba\)。在第一个计算中，代表项 \(a\) 消失；在第二个计算中，多出了一个来自另一因子的代表项。这正是证明中两种前缀操作的具体情形。对任意输入字重复这些步骤即可化简；其有效计算在\subsecref{b1:subsec:C03:02}中继续说明。

\ProseParagraphBreak
上例的公共元素 \(z\) 恰在中心，不能用它代替一般的前缀分解。再取两份 \(S_3\)：
\[
A=\langle r,s\mid r^3=s^2=1,\ srs=r^{-1}\rangle,
\qquad
B=\langle t,u\mid t^3=u^2=1,\ utu=t^{-1}\rangle.
\]
沿同构 \(\langle s\rangle\cong\langle u\rangle\) 合并后，把两者在 \(A*_{C_2}B\) 中的公共像仍记为 \(s\)。可取右陪集代表集 \(T_A=\{1,r,r^2\}\)、\(T_B=\{1,t,t^2\}\)。在标准串 \((s;t)\) 上左乘 \(r\) 时，关系 \(srs=r^{-1}\) 给出 \(rs=sr^2\)，所以
\[
L_r(s;t)=(s;r^2,t).
\]
若把 \(r\) 与 \(s\) 当作可交换元素，便会误得 \((s;r,t)\)；两个串的代表项不同，按标准形定理表示不同元素。一般而言，把公共子群系数移回左端需要重新作 \(ct\)-分解，代表项可能随之改变，不能将这一步理解为把所有 \(C\) 元素与代表逐一交换。
